% SPDX-License-Identifier: LPPL-1.3c
% Copyright (C) 2026 Christophe Jorssen
% Author and Current Maintainer: Christophe Jorssen <christophe.jorssen@gmail.com>
% LPPL maintenance status: maintained
% This file is part of luafemm. See LICENSE and LICENSES.md for its terms.
\section{Native FEMM air-gap elements}
\label{sec:native-gaps}

A Fourier annulus and a native FEMM air-gap element (AGE) are different numerical
models. Choose the native formulation when the saved answer must carry FEMM's
own gap records. Choose |analytical| when the exact radial Fourier formulation
or sector reduction is wanted. Neither option silently converts the other's
solved field.

\begin{key}{/femm/sliding band/method=femm4}
Select the native cubic interpolation for a full-circle concentric planar gap.
The ordinary initial choice remains |p1|. The inner and outer traces must have
the same number of uniformly spaced angular nodes; increase |segments| if bulk
refinement subdivides a trace. The total native trace operator is limited to
2048 nodes and 2048 total cells (at least three per gap). |layers| is ignored (one display strip); |modes| must be zero.
\end{key}
\begin{key}{/femm/sliding band/method=femm5}
Select the native quintic interpolation, with the same geometry and size limits.
This is the appropriate explicit choice for the supported FEMM 5 solver.
It does not change the default |p1| method. Both native choices support static
and linear harmonic positions, including the supported material laws in the
bulk. They require zero rotor offsets and one full-circle sector.
\end{key}

\luafemmexamplebegin{native-gap}
\begin{codeexample}[width=0pt]
\tikzset{femm/materials/gap iron/.style={mu r=100}}
\begin{tikzpicture}[
  femm/problem={xmin=-8,xmax=8,ymin=-8,ymax=8,mesh size=1.2,depth=10},
  femm/sliding band={name=native-ring,method=femm5,
    inner radius=3,outer radius=3.6,segments=64},
  scale=2.5,baseline=(current bounding box.center)]
  \draw[femm/region={material=gap iron}]
    (0,0) ellipse[x radius=2.6,y radius=1.2];
  \draw[femm/region={material=copper,current density=1e7}]
    (-6,-1) rectangle (-5,1);
  \draw[femm/region={material=copper,current density=-1e7}]
    (5,-1) rectangle (6,1);
\end{tikzpicture}
\begin{tikzpicture}[
  femm/problem={sliding band=native-ring,rotor angle=25},
  scale=2.5,baseline=(current bounding box.center)]
  \pic {femm geometry};
  \pic[femm/lines=13] {femm field};
  \femmexport[file=\jobname-native-gap.fem,format=fem]
  \femmexport[file=\jobname-native-gap.ans,format=ans]
\end{tikzpicture}
\end{codeexample}

The first picture declares the reference geometry. The second solves at
$25^\circ$ and exports that physical position. The exported native inner and
outer angle offsets are zero: the coordinates already contain the rotation.
Its circular gap boundaries remain arcs in the problem file. Arbitrary TikZ
curves elsewhere retain the usual controlled polygonal approximation.

\Needspace{6\baselineskip}
\begin{key}{/femm/import/age interpolation=\meta{choice} (initially empty)}
For an ordinary planar AGE answer, explicitly choose |femm4| (cubic) or |femm5|
(quintic). The file format does not reliably identify the interpolation basis;
the optional raw-ring trailer is not a basis tag. A missing choice is an error.
This key belongs to |femm/import|, applies to |format=ans|, and has no units.
It neither changes material interpretation nor converts an answer to another
basis. Transient and previous-solution AGE answers remain unsupported.
\end{key}
\begin{codeexample}[]
\begin{tikzpicture}[
  femm/import={file=\jobname-native-gap.ans,
    age interpolation=femm5},scale=2.5]
  \pic {femm geometry};
  \pic[femm/lines=13] {femm field};
\end{tikzpicture}
\end{codeexample}
\luafemmexampleend{native-gap}

An imported answer owns its mesh and potentials; importing it does not perform
a new solve. The embedded materials and circuits must satisfy the ordinary
interchange limits. Periodic nodal values are checked, and native AGE names,
indices, dimensions and optional raw rings are validated. Re-exporting an
unchanged native snapshot preserves its records. General AGE problem-file
import for remeshing is not supported: import the native solved |.ans| instead.
Native answer generation currently covers the full-circle templates above;
sector AGE answers can be read, but sector generation uses the Fourier method.

\subsection{Which field and which mechanical result?}
The native stiffness is assembled as $T^\mathsf{T}K_\square T$: $T$ interpolates
four corner potentials from five nodes on each trace. Cubic Hermite interpolation
uses central slopes; quintic interpolation also uses central second differences.
The rectangular energy combines the triangular and bilinear forms with weights
$2/3$ and $1/3$. This factorisation is checked independently against the native
expanded matrices for both editions.

Gap samples on the mean radius use the native postprocessor's Fourier transform
of cell-centre $B_r,B_\theta$. Away from that circle, luafemm continues those
harmonics radially with Laplace's equation. This continuation is distinct from
bulk P1 interpolation; do not use it as an independent discretisation certificate.
The native formulation approximates a thin annulus by an unrolled strip.
Check both angular refinement and gap thickness relative to radius.

|\femmmeasure[quantity=force,method=band,...]| integrates the gap traction;
|origin x| and |origin y| determine the torque origin. It also accepts a native
AGE snapshot. Energy uses the native stiffness, while torque uses the native
midgap field; their virtual-work agreement is a convergence check, not an exact
finite-mesh identity. In an imported file with several AGE blocks, the band
measurement sums their transmitted tractions. It is not a selection of one rotor.
The gap views sample the reconstructed potential, controlled by
|contour subdivisions|; separate bulk/gap pieces are not forcibly joined.
Native bulk coordinates follow the file's body-frame convention. Nonzero saved
native angle offsets are not interpreted as a rigid rotation of every bulk part.

Regional areas and mean fields integrate the available volume triangles. An
imported native answer has no such triangles inside its AGE blocks, whereas a
generated template has a display strip. Consequently these regional quantities
are not an interchange invariant. Whole-model energy includes the native AGE
energy in both cases; a partial bulk selection in an imported answer does not
select an AGE block. Use gap samples and band traction for gap-specific results.

The common AGE records are readable by FEMM 4 and FEMM 5, but their identical
syntax does not make the cubic and quintic operators interchangeable. Open a
|method=femm4| answer with the cubic solver/postprocessor, and a |method=femm5|
answer with the quintic one. The latter also carries |[AgeRings]| for native
re-pairing; omitting that trailer does not turn a quintic field into a cubic one.
